% GENERATED_BY: oc_core_monograph_translation_draft_pipeline_v1.py
% AI_ASSISTED_TRANSLATION_DRAFT: TRUE
% THEOREM_REVIEW_STATUS: PENDING
% THEOREM_REVIEW_MODE: FORMAL_AI_GATE__CODEX_CLI_CHATGPT
% THEOREM_REVIEW_REVIEWER_ID:
% THEOREM_REVIEW_AUTHORITY_CLASS:
% THEOREM_REVIEWED_AT:
% SOURCE_LANGUAGE: EN
% TARGET_LANGUAGE: RU
% SOURCE_EN_REF: content/k_levels/k4.tex
% ================================================================
% ==== FILE: content/k_levels/k4.tex
% ================================================================

% ==============================
%  Ontology of Continua — Core
%  K-level Module: K4
%  File: content/k_levels/k4.tex
%  Status: FULLY DEFINED
% ==============================

\subsubsection{\texorpdfstring{$K_4$}{K_4} Обзор}
\label{sec:k4-overview}

$K_4$ — уровень, на котором впервые становится возможной
\emph{биологическая} структура. Ключевое событие — появление
\emph{полу-устойчивой мембраны}, которая создаёт новое онтологическое
разделение:
\[
\text{внутри} \;|\; \text{снаружи}.
\]

Это действительно новая ось континуума, порождающая:
\begin{itemize}
    \item ограниченную химическую внутреннюю область,
    \item селективные режимы переноса,
    \item пространственно структурированные градиенты,
    \item прото-гомеостаз,
    \item первые биологические циклы возбуждения, выживания и восстановления.
\end{itemize}

Переход $K_3 \to K_4$ соответствует формированию протоклетки: RAF-
сети в полупроницаемой мембране, чья стабильность определяется
пороговыми условиями $\Theta_{\mathrm{mem}}$,
$\Theta_{\mathrm{perm}}$, $\Theta_{\mathrm{osm}}$ и
$\Theta_{\mathrm{curv}}$.

$K_4$ — минимальный биологический континуум: наиболее простая форма
организованной и устойчивой протоклеточной жизни.

% ================================================================
\subsubsection{Пространство состояний \texorpdfstring{$\Omega(K_4)$}{\Omega(K_4)}}

Пространство состояний включает:
\[
\Omega(K_4) =
\{(\mathbf{c}_{\mathrm{in}},\mathbf{c}_{\mathrm{out}},
  \partial\Omega,\; \mathbf{g},\; \Gamma_{\mathrm{RAF}},
  \mathbf{p})\},
\]
где:
\begin{itemize}
    \item $\mathbf{c}_{\mathrm{in}}$: концентрации внутри мембраны,
    \item $\mathbf{c}_{\mathrm{out}}$: внешние концентрации,
    \item $\partial\Omega$: граница мембраны с участками (patch-structure),
    \item $\mathbf{g}$: градиенты через мембрану (осмотический, ионный,
          pH, потенциал),
    \item $\Gamma_{\mathrm{RAF}}$: внутренняя каталитическая сеть,
    \item $\mathbf{p}$: параметры проницаемости и кривизны для каждого
          участка границы.
\end{itemize}

Пространство состояний имеет более высокую размерность по сравнению с $K_3$
благодаря мембранной топологии и степеням свободы по участкам.

Допустимые состояния должны удовлетворять:
\begin{itemize}
    \item когезии мембраны,
    \item контролируемой проницаемости,
    \item некритичным градиентам (без разрушительного эффекта),
    \item жизнеспособности RAF-сети внутри.
\end{itemize}

% ================================================================
\subsubsection{\texorpdfstring{$\partial\Omega(K_4)$}{\partial\Omega(K_4)} (граница)}

Мембрана — это \emph{динамическая граница}:
\[
\partial\Omega(K_4) = \bigcup_{i=1}^{N_{\mathrm{patch}}}
    \partial\Omega_i.
\]

Каждый участок $\partial\Omega_i$ имеет:
\begin{itemize}
    \item локальную кривизну $K_i$,
    \item локальную проницаемость $p_i$,
    \item локальное натяжение $T_i$,
    \item локальные ограничения градиентов (осмотический, ионный, pH),
    \item фазовое состояние (L\textsubscript{α}, L\textsubscript{o},
          L\textsubscript{β}, пористое).
\end{itemize}

Участок разрушается, если:
\[
T_i > \Theta_{\mathrm{mem},i}
\quad \text{или} \quad
\Delta \Pi_i > \Theta_{\mathrm{osm},i}
\quad \text{или} \quad
K_i > \Theta_{\mathrm{curv},i}.
\]

Потеря любого критического участка уничтожает $\Omega(K_4)$.

% ================================================================
\subsubsection{Оси \texorpdfstring{$A(K_4)$}{A(K_4)}}

$K_4$ поддерживает по крайней мере следующие оси:
\[
A(K_4) = \{ A_{\mathrm{comp}}, A_{\mathrm{boundary}},
           A_{\mathrm{curv}}, A_{\mathrm{perm}}, A_{\mathrm{grad}}\}.
\]

Интерпретация:
\begin{itemize}
    \item $A_{\mathrm{comp}}$: различие компартмента внутри/снаружи,
    \item $A_{\mathrm{boundary}}$: степени свободы мембраны,
    \item $A_{\mathrm{curv}}$: состояния кривизны участков границы,
    \item $A_{\mathrm{perm}}$: режимы переменной проницаемости,
    \item $A_{\mathrm{grad}}$: химические/электрические/ионные градиенты.
\end{itemize}

Эти оси \emph{не} сводятся к химическим осям $K_3$; каждая
отражает новые структурные различия биологической организации.

% ================================================================
\subsubsection{Потенциалы \texorpdfstring{$P(K_4)$}{P(K_4)}}

Включённые потенциалы:
\[
P(K_4) =
\{ \mu_i^{\mathrm{in}},\; \mu_i^{\mathrm{out}},
   P_{\mathrm{grad}},\;
   P_{\mathrm{osm}},\;
   P_{\mathrm{curv}},\;
   P_{\mathrm{mem}},\;
   P_{\mathrm{pump}} \}.
\]

Значения:
\begin{itemize}
    \item внутренние/внешние химические потенциалы,
    \item потенциалы градиентов, задающих селективный перенос,
    \item потенциалы осмотического давления,
    \item энергия кривизны участков мембраны,
    \item потенциалы натяжения мембраны,
    \item энергия, запасённая в примитивных насосах (если они есть).
\end{itemize}

Эти потенциалы регулируют выживаемость: слишком большой градиент или
кривизна разрушают систему.

% ================================================================
\subsubsection{Пороги \texorpdfstring{$\Theta(K_4)$}{\Theta(K_4)}}

Пороги определяют устойчивость мембраны и протоклетки:
\[
\Theta(K_4) =
\{\Theta_{\mathrm{mem}},\;
  \Theta_{\mathrm{perm}},\;
  \Theta_{\mathrm{osm}},\;
  \Theta_{\mathrm{curv}},\;
  \Theta_{\mathrm{pH}},\;
  \Theta_{\mathrm{grad}},\;
  \Theta_{\mathrm{RAF-survival}}\}.
\]

Ключевые интерпретации:
\begin{itemize}
    \item $\Theta_{\mathrm{mem}}$: максимальное натяжение мембраны перед
          разрывом,
    \item $\Theta_{\mathrm{perm}}$: пределы пассивной проницаемости,
    \item $\Theta_{\mathrm{osm}}$: осмотическая выносливость перед разрывом,
    \item $\Theta_{\mathrm{curv}}$: порог нестабильности кривизны,
    \item $\Theta_{\mathrm{pH}}$: допустимый внутренний диапазон pH,
    \item $\Theta_{\mathrm{grad}}$: устойчивые градиенты без коллапса,
    \item $\Theta_{\mathrm{RAF-survival}}$: минимальная каталитическая активность
          для гомеостатических циклов.
\end{itemize}

Эти пороги совместно определяют границу допустимых состояний протоклетки.

% ================================================================
\subsubsection{Потоки \texorpdfstring{$J(K_4)$}{J(K_4)}}

Потоки в $K_4$ включают:
\begin{itemize}
    \item трансмембранный транспорт:
          \[
          J_{\mathrm{ion}},\; J_{\mathrm{water}},\; J_{\mathrm{solute}},
          \]
    \item пассивные утечки потоков,
    \item потоки, обусловленные кривизной, и перестройка участков,
    \item внутренние RAF-реактивные потоки,
    \item насоса-зависимые активные потоки (если есть примитивные насосы),
    \item циклы осмотического набухания–релаксации.
\end{itemize}

Для устойчивости потоков требуется контроль:
\[
\Delta V,\; \Delta \Pi,\; \Delta c_i,\; \Delta \mathrm{pH}.
\]

% ================================================================
\subsubsection{Циклы \texorpdfstring{$C(K_4)$}{C(K_4)}}

Биологические циклы появляются первыми:
\[
C(K_4) =
\{ C_{\mathrm{survival}},\;
   C_{\mathrm{osmotic}},\;
   C_{\mathrm{pH-recovery}},\;
   C_{\mathrm{leak-pump}},\;
   C_{\mathrm{RAF-stability}} \}.
\]

Значения:
\begin{itemize}
    \item $C_{\mathrm{survival}}$: минимальный гомеостатический цикл, сохраняющий
          целостность мембраны,
    \item $C_{\mathrm{osmotic}}$: повторяющиеся циклы набухания–релаксации,
    \item $C_{\mathrm{pH-recovery}}$: буферизация и реакционные потоки,
    \item $C_{\mathrm{leak-pump}}$: баланс пассивных утечек и активного
          или полуактивного переноса,
    \item $C_{\mathrm{RAF-stability}}$: внутренние динамические сети,
          поддерживающие стабильность границы.
\end{itemize}

Циклы определяют биологическую направленность времени и устойчивость
протоклетки.

% ================================================================
\subsubsection{Время \texorpdfstring{$\tau(K_4)$}{\tau(K_4)}}

Время существует, когда:
\[
\Pi(C_{\mathrm{survival}}) > \Theta_{\mathrm{time}}
\quad\text{и}\quad
\Pi(C_{\mathrm{RAF-stability}}) > \Theta_{\mathrm{time}}.
\]

Следовательно:
\begin{itemize}
    \item биологическое время возникает из циклов выживания,
    \item время коллапсирует, когда мембрана не может поддерживать
          устойчивую работу циклов,
    \item временной порядок более выражен, чем в $K_3$, из-за обратных связей
          границы.
\end{itemize}

% ================================================================
\subsubsection{Континуальность \texorpdfstring{$k(K_4)$}{k(K_4)}}

Применяя общую формулу:
\[
k(K_4) =
\frac{\mu(\Omega(K_4))}{\mu(S(K_4))}
\cdot
\frac{|A(K_4)|}{|A(M_4)|}
\cdot
\frac{\sum \mathrm{Stab}(C)}{\sum \mathrm{MaxStab}(C)}
\cdot
(1 - \frac{\sigma(J)}{\sigma_{\max}})
\cdot
(1 - \frac{T}{\Theta_{\mathrm{stab}}})_+.
\]

Следствия:
\begin{itemize}
    \item стабильность мембраны увеличивает $k(K_4)$,
    \item осмотические и кривизные флуктуации уменьшают $k(K_4)$,
    \item потеря целостности участков задаёт $k(K_4)=0$,
    \item успешная организация протоклетки соответствует
          $k(K_4)>0$ на длительный период.
\end{itemize}

% ================================================================
\subsubsection{Структурное напряжение \texorpdfstring{$T(K_4)$}{T(K_4)}}

Источники:
\begin{itemize}
    \item несовместимые осмотические градиенты,
    \item чрезмерное кривизное напряжение на участках границы,
    \item pH-стресс,
    \item неконтролируемая проницаемость,
    \item дисбаланс между утечкой и насосными потоками,
    \item недостаточная поддержка RAF.
\end{itemize}

Структурный коллапс происходит, если:
\[
T(K_4) > \Theta_{\mathrm{mem}}.
\]

% ================================================================
\subsubsection{Энергия \texorpdfstring{$E(K_4)$}{E(K_4)}}

Энергетический функционал включает:
\[
E(K_4) =
E_{\mathrm{chem}}^{\mathrm{in}} + E_{\mathrm{chem}}^{\mathrm{out}}
+ E_{\mathrm{osm}} + E_{\mathrm{curv}}
+ E_{\mathrm{grad}} + E_{\mathrm{mem}}
+ E_{\mathrm{RAF}}.
\]

Компоненты:
\begin{itemize}
    \item внутренняя и внешняя химические энергии,
    \item энергия осмотического давления,
    \item энергия кривизны мембраны,
    \item энергия градиента,
    \item энергия натяжения мембраны,
    \item каталитическая энергия сети RAF.
\end{itemize}

$E(K_4)$ определяет стабильность и динамику протоклетки.

% ================================================================
\subsubsection{Операторы на \texorpdfstring{$K_4$}{K_4} (\texorpdfstring{$\Psi$}{\Psi}, \texorpdfstring{$\Phi$}{\Phi}, \texorpdfstring{$\Lambda$}{\Lambda}, \texorpdfstring{$U$}{U}, \texorpdfstring{$\Chi$}{\Chi})

\begin{itemize}
    \item $\Psi_{4\to5}$ — возникновение электрической возбудимости и
          примитивных ионных каналов (переход к $K_5$),
    \item $\Phi$ — эволюция ограничений мембраны и разрешённых химических
          состояний,
    \item $\Lambda$ — структурная связь между мембраной и наборами RAF,
    \item $U$ — унификация подсистем градиента и мембраны,
    \item $\Chi$ — ветвление в альтернативные типы протоклеток (разные
          химии мембраны или режимы переноса).
\end{itemize}

% ================================================================
\subsubsection{Процессы на \texorpdfstring{$K_4$}{K_4}}

Характерные процессы:
\begin{itemize}
    \item образование и закрытие везикул,
    \item перестройки, обусловленные кривизной,
    \item осмотическое набухание и релаксация,
    \item циклы буферизации pH,
    \item динамика баланса утечки и насоса,
    \item протометрические обратные петли,
    \item возникновение ранних ионных каналов,
    \item стабилизация электрической оси (pre-$K_5$).
\end{itemize}

% ================================================================
\subsubsection{Прогнозы для \texorpdfstring{$K_4$}{K_4}}

\begin{enumerate}
    \item Для устойчивости мембраны требуют пороги кривизны и осмоса в
          узких границах.
    \item Распределение натяжения по участкам предсказывает локальные
          режимы мерцания.
    \item Циклы восстановления pH необходимы для продолжения протоклетки.
    \item Сети RAF по отдельности не поддерживают организацию без
          устойчивой стабильности границы.
    \item Электрические градиенты начинают проявляться до полного $K_5$.
    \item Переход к $K_5$ возможен только после стабилизации ионных
          каналов.
\end{enumerate}

% ================================================================
\subsubsection{Эксперименты для \texorpdfstring{$K_4$}{K_4}}

Эмпирические аналоги включают:
\begin{itemize}
    \item эксперименты с везикулами жирных кислот,
    \item тесты протоклеточного осмотического разрыва,
    \item картирование на уровне участков флуоресценцией мерцания мембраны,
    \item вставка ионных каналов в простые везикулы,
    \item сети RAF, заключённые во везикулы,
    \item динамика пH-буферных везикул.
\end{itemize}

Эти тесты проверяют пороги $\Theta_{\mathrm{mem}}$,
$\Theta_{\mathrm{curv}}$, $\Theta_{\mathrm{osm}}$,
$\Theta_{\mathrm{pH}}$.

% ================================================================
\subsubsection{Коллапс и смерть \texorpdfstring{$K_4$}{K_4}}

Коллапс происходит, когда:
\[
\Omega(K_4) = \varnothing.
\]

Причины:
\begin{itemize}
    \item разрыв мембраны,
    \item экстремальный осмотический дисбаланс,
    \item разрыв, вызванный кривизной,
    \item катастрофический коллапс pH,
    \item неконтролируемая утечка,
    \item отказ циклов выживания.
\end{itemize}

После коллапса переход в $K_5$ невозможен.

% ================================================================
\subsubsection{Фальсифицируемость \texorpdfstring{$K_4$}{K_4}}

Прогнозы, которые могут быть экспериментально проверены:
\begin{enumerate}
    \item существование строгих мембранных порогов для выживания,
    \item универсальная форма осмотических кривых коллапса,
    \item необходимость RAF-поддерживаемых буферных циклов,
    \item невозможность устойчивой организации протоклетки без контроля
          кривизны на уровне участков,
    \item раннее появление прото-электрических градиентов,
    \item требование минимального баланса утечка–накачка.
\end{enumerate}

Отрицание этих прогнозов опровергнет определение $K_4$ в текущем виде.

% ================================================================
\subsubsection{Разветвление / Онтологическое положение \texorpdfstring{$K_4$}{K_4}}

Типы ветвей:
\begin{itemize}
    \item различные химии мембраны (жирные кислоты, фосфолипиды,
          смешанные),
    \item различные режимы проницаемости,
    \item альтернативные ядра RAF, совместимые с мембранной динамикой,
    \item расходящиеся протометаболические стратегии.
\end{itemize}

Онтологическое положение:
\[
K_3 \to K_4 \to K_5,
\]
где $K_4$ — первый биологический континуум.

% ================================================================
\subsubsection{Связь с M-пространствами}

$M_4$ задаёт:
\begin{itemize}
    \item допустимые мембранные материалы,
    \item разрешённые диапазоны кривизны и проницаемости,
    \item средовые режимы (осмотический, ионный, pH),
    \item допустимые процессы переноса,
    \item ограничения на формирование $\partial\Omega$ как устойчивой границы.
\end{itemize}

Совместимость $K_4$ с $M_4$ определяет жизнеспособность протоклетки.
 \text{внутри} \;|\; \text{снаружи}.
